% !TEX program = xelatex
% 示例论文：用于验证 well-posed 的整条管线能否跑通。
% 声明：这不是一篇真实发表的论文，是为了测试而写的短文；
% 其中的数学结论是正确的、可手工验证的，但没有任何研究价值。
\documentclass[11pt]{article}
\usepackage[a4paper,margin=2.6cm]{geometry}
\usepackage{amsmath,amssymb,amsthm}
\usepackage{booktabs}
\usepackage[colorlinks=true,linkcolor=blue,citecolor=blue]{hyperref}

\newtheorem{theorem}{Theorem}
\newtheorem{proposition}{Proposition}
\newtheorem{lemma}{Lemma}
\newtheorem{remark}{Remark}
\theoremstyle{definition}
\newtheorem{definition}{Definition}

\title{Exponential Decay for a Damped Wave Equation:\\ A Minimal Test Fixture}
\author{Test Fixture$^{\dagger}$\\
\small $^{\dagger}$This is \emph{not} a published paper. It exists only to exercise the\\
\small well-posed pipeline end to end. The mathematics below is correct but elementary.}
\date{2026}

\begin{document}
\maketitle

\begin{abstract}
We consider the damped wave equation $u_{tt}+u_{t}-u_{xx}=0$ on a bounded interval
with homogeneous Dirichlet boundary conditions. We give a self-contained proof that the
energy decays exponentially, with an explicit rate obtained from a multiplier identity.
The note is written as a test fixture: it is short, it contains exactly one skipped step
(so that a reader may practice filling it in), and its numbers can be checked by hand.
\end{abstract}

\section{Introduction and main result}

Let $L>0$ and $\Omega=(0,L)$. Consider
\begin{equation}
u_{tt}+u_{t}-u_{xx}=0,\qquad x\in(0,L),\ t>0,
\label{eq:model}
\end{equation}
with Dirichlet boundary conditions
\begin{equation}
u(0,t)=u(L,t)=0,\qquad t\ge 0,
\label{eq:bc}
\end{equation}
and initial data
\begin{equation}
u(x,0)=u_{0}(x),\qquad u_{t}(x,0)=u_{1}(x).
\label{eq:ic}
\end{equation}

The energy of a solution is
\begin{equation}
E(t)=\frac12\int_{0}^{L}\bigl(u_{t}^{2}+u_{x}^{2}\bigr)\,dx .
\label{eq:energy}
\end{equation}

\begin{proposition}[Energy identity]
\label{prop:energy}
Let $u$ be a sufficiently smooth solution of \eqref{eq:model}--\eqref{eq:bc}.
Then
\begin{equation}
\frac{d}{dt}E(t)=-\int_{0}^{L}u_{t}^{2}\,dx\le 0 .
\label{eq:dissipation}
\end{equation}
\end{proposition}

\begin{proof}
Multiply \eqref{eq:model} by $u_{t}$ and integrate over $(0,L)$:
\[
\int_{0}^{L}u_{tt}u_{t}\,dx+\int_{0}^{L}u_{t}^{2}\,dx-\int_{0}^{L}u_{xx}u_{t}\,dx=0 .
\]
The first term is $\frac12\frac{d}{dt}\int u_{t}^{2}$, and integration by parts in the
third term, together with \eqref{eq:bc}, gives
$-\int u_{xx}u_{t}=\int u_{x}u_{xt}=\frac12\frac{d}{dt}\int u_{x}^{2}$.
Collecting terms yields \eqref{eq:dissipation}.
\end{proof}

Proposition~\ref{prop:energy} shows that $E$ is nonincreasing, but it does not by itself
give a rate. The main result of this note is the following.

\begin{theorem}[Exponential decay]
\label{thm:decay}
Let $u$ solve \eqref{eq:model}--\eqref{eq:ic} with
$(u_{0},u_{1})\in H^{2}(\Omega)\cap H^{1}_{0}(\Omega)\times H^{1}_{0}(\Omega)$.
Then there exist constants $C_{1}\ge 1$ and $\mu>0$, depending only on $L$, such that
\begin{equation}
E(t)\le C_{1}\,e^{-\mu t}E(0),\qquad t\ge 0 .
\label{eq:decay}
\end{equation}
\end{theorem}

\begin{remark}
The proof below is the classical multiplier argument. It is included in full because the
constants are explicit and can be checked with a pencil; readers are invited to verify
that the choice \eqref{eq:eps} below is legitimate.
\end{remark}

\section{Proof of Theorem~\ref{thm:decay}}

Define the auxiliary functional
\begin{equation}
F(t)=\int_{0}^{L}u\,u_{t}\,dx+\frac12\int_{0}^{L}u^{2}\,dx .
\label{eq:F}
\end{equation}

\begin{lemma}
\label{lem:F}
Along solutions of \eqref{eq:model}--\eqref{eq:bc},
\begin{equation}
F'(t)=\int_{0}^{L}u_{t}^{2}\,dx-\int_{0}^{L}u_{x}^{2}\,dx .
\label{eq:Fprime}
\end{equation}
\end{lemma}

\begin{proof}
Differentiating \eqref{eq:F} and using \eqref{eq:model} in the form
$u_{tt}=u_{xx}-u_{t}$,
\begin{align*}
F'(t)&=\int_{0}^{L}\bigl(u_{t}^{2}+u\,u_{tt}\bigr)dx+\int_{0}^{L}u\,u_{t}\,dx\\
&=\int_{0}^{L}u_{t}^{2}\,dx+\int_{0}^{L}u\,(u_{xx}-u_{t})\,dx+\int_{0}^{L}u\,u_{t}\,dx\\
&=\int_{0}^{L}u_{t}^{2}\,dx-\int_{0}^{L}u_{x}^{2}\,dx ,
\end{align*}
where the boundary term $\bigl[u\,u_{x}\bigr]_{0}^{L}$ vanishes by \eqref{eq:bc}
and the two terms involving $u\,u_{t}$ cancel.
\end{proof}

Now set, for $\varepsilon>0$ to be chosen,
\begin{equation}
E_{\varepsilon}(t)=E(t)+\varepsilon F(t).
\label{eq:Eeps}
\end{equation}

Combining \eqref{eq:dissipation} and \eqref{eq:Fprime},
\begin{equation}
E_{\varepsilon}'(t)=-(1-\varepsilon)\int_{0}^{L}u_{t}^{2}\,dx-\varepsilon\int_{0}^{L}u_{x}^{2}\,dx .
\label{eq:Eepsprime}
\end{equation}

It remains to compare $E_{\varepsilon}$ with $E$. By the Cauchy--Schwarz inequality and
the Poincaré inequality one obtains
\begin{equation}
|F(t)|\le C_{0}E(t),\qquad C_{0}=C_{0}(L),
\label{eq:equiv}
\end{equation}
so that for $\varepsilon$ small enough,
\begin{equation}
(1-\varepsilon C_{0})E(t)\le E_{\varepsilon}(t)\le(1+\varepsilon C_{0})E(t).
\label{eq:compare}
\end{equation}

Choose
\begin{equation}
0<\varepsilon<\frac{1}{2C_{0}},\qquad \varepsilon<\frac12 .
\label{eq:eps}
\end{equation}
Then \eqref{eq:Eepsprime} gives, for some $c=c(\varepsilon)>0$,
$E_{\varepsilon}'\le-cE$, and by \eqref{eq:compare},
$E_{\varepsilon}'\le-\frac{c}{1+\varepsilon C_{0}}E_{\varepsilon}$.
Gronwall's inequality yields
\[
E_{\varepsilon}(t)\le e^{-\mu t}E_{\varepsilon}(0),\qquad
\mu=\frac{c}{1+\varepsilon C_{0}},
\]
and applying \eqref{eq:compare} once more gives \eqref{eq:decay} with
$C_{1}=(1+\varepsilon C_{0})/(1-\varepsilon C_{0})$.

\begin{remark}
The estimate \eqref{eq:equiv} is the only place where the boundary conditions enter the
argument beyond \eqref{eq:dissipation}; without \eqref{eq:bc} the Poincaré inequality is
unavailable and $F$ cannot be dominated by $E$.
\end{remark}

\section{Conclusion}

The proof uses three ingredients: an energy identity
(Proposition~\ref{prop:energy}), a multiplier identity producing a functional whose
derivative has the opposite sign (Lemma~\ref{lem:F}), and an equivalence estimate that
turns the modified energy into the original one. This is the same scheme used for
variable-coefficient and higher-dimensional versions of \eqref{eq:model}.

\begin{thebibliography}{9}
\bibitem{Komornik} V.~Komornik, \emph{Exact Controllability and Stabilization: The
Multiplier Method}, Masson/Wiley, 1994.
\bibitem{Lions} J.-L. Lions, \emph{Contrôlabilité exacte, perturbations et stabilisation
de systèmes distribués}, Masson, 1988.
\bibitem{Zuazua} E.~Zuazua, Exponential decay for the semilinear wave equation with
locally distributed damping, \emph{Comm. PDE} 15 (1990), 205--235.
\end{thebibliography}

\end{document}
